\section{Related Work}
\label{sec:intro-related-work}

The work presented in this manuscript lies at the confluence of three research traditions that have largely developed on their own terms: deductive verification frameworks, which reduce the correctness of software to logical proof obligations; automated theorem proving, and satisfiability modulo theories (SMT) solving in particular; and the mechanization of languages and their metatheory inside proof assistants.
This section retraces the relevant history of each tradition and then reviews prior work on the problem at the heart of this thesis: translating proof obligations rooted in set theory into the language of automated provers, and the degree of trust such translations deserve.

\subsection{Proof obligations and the cost of proof}
\label{subsec:rw-proof-obligations}

Deductive program verification goes back to the assertion method of Floyd~\cite{floyd1967meanings} and Hoare's axiomatic semantics~\cite{hoare1969axiomatic}, later recast by Dijkstra's weakest-precondition calculus~\cite{dijkstra1975guarded} into a systematic discipline: the correctness of a program with respect to its specification is reduced to a finite collection of logical formulas, the \emph{verification conditions} or \emph{proof obligations}, whose validity implies that of the program.
Under this reduction, the practicality of formal development is governed by a single factor: the effort needed to discharge these formulas.

The B method~\cite{abrial1996bbook} is a prominent industrial embodiment of this idea: every step of a B development generates proof obligations expressed in the set-theoretic language of the method (\Cref{sec:b-method} presents it in detail).
The method earned its reputation with the automatic train protection software of the Paris M\'etro Line~14~\cite{behm1999meteor}, deployed without a single defect traced to the formally developed code, and it has since sustained a continuous industrial practice in railway signaling and other safety-critical domains~\cite{butler2020first}.
Its descendant Event-B~\cite{abrial2010modeling} streamlines the notation and generalizes refinement, and is supported by the extensible Rodin platform~\cite{abrial2010rodin}.

In both settings, proof is the dominant cost.
Atelier~B, the industrial development environment for B~\cite{lecomte2024proving}, discharges obligations with its built-in automatic provers, the predicate prover (PP) and the mono-lemma prover (ML); every obligation they miss falls back to interactive proof or to user-added rules whose validity must in turn be justified.
On industrial projects, which routinely generate tens of thousands of obligations, even a residual few percent of unproved obligations translates into weeks of expert work.
This economic pressure is the original motivation for connecting B tools to external automated provers~\cite{deharbe2010automatic,mentre2012discharging}---and each such connection immediately raises the question that runs through this entire survey: the connection is a \emph{translation} between logics, and every guarantee obtained downstream is conditioned on its soundness.

\subsection{From resolution to satisfiability modulo theories}
\label{subsec:rw-atp-smt}

Automated deduction predates mechanized verification.
Robinson's resolution principle~\cite{robinson1965resolution} gave first-order logic a single machine-oriented inference rule and founded a lineage of saturation-based provers that continues with Vampire~\cite{kovacs2013vampire}, E~\cite{schulz2013e}, and Zipperposition~\cite{bentkamp2023higherorder}.
Set theory was an early and revealing stress test for this technology: Quaife's semi-automated development of von Neumann--Bernays--G\"odel set theory with the Otter prover~\cite{quaife1992nbg} demonstrated both that resolution could in principle carry elementary set theory and how laborious axiomatic reasoning about membership remains for general-purpose provers---a difficulty that every encoding surveyed below has had to confront.

A second lineage grew out of propositional logic.
The procedures of Davis, Putnam, Logemann, and Loveland~\cite{davis1960computing,davis1962machine} evolved, with conflict-driven clause learning (CDCL)~\cite{marquessilva1999grasp,biere2021handbook}, into SAT solvers of remarkable scalability.
Program checking, however, needs more than propositional reasoning: it needs equality, arithmetic, and data structures.
Nelson and Oppen showed how decision procedures for such theories cooperate~\cite{nelson1979simplification}, and their approach matured into Simplify~\cite{detlefs2005simplify}, the prover behind the extended static checker ESC/Java~\cite{flanagan2002escjava} and the direct ancestor of modern SMT solvers.
The synthesis of both lineages, formalized as DPLL(\(T\))~\cite{nieuwenhuis2006dpllt}, combines CDCL-style propositional search with dedicated theory solvers and defines the architecture of the field's flagship systems: Z3~\cite{demoura2008z3}, CVC4 and its successor cvc5~\cite{barrett2011cvc4,barbosa2022cvc5}, veriT~\cite{bouton2009verit}, and Alt-Ergo~\cite{conchon2018altergo}, the last natively supporting the polymorphic first-order logic favored by deductive verification platforms.
The SMT-LIB initiative~\cite{barrett2010smtlib} standardized the input language and, together with the annual solver competition, turned SMT solving into a reliable off-the-shelf technology~\cite{demoura2011smt}.
Its expressiveness has kept growing in directions relevant to this thesis: a decision procedure for finite sets with cardinality~\cite{bansal2016sets}, pragmatic extensions of SMT solvers to higher-order logic~\cite{barbosa2019hosmt}, and, with version~2.7 of the standard, lambda abstractions in the language itself~\cite{smtlib2.7}.
\Cref{sec:ar-smt} presents the technical background on SMT solving assumed in the rest of the manuscript.

\subsection{Intermediate verification languages}
\label{subsec:rw-ivl}

With reliable automated back-ends available, the verification community converged on a common architecture: a front-end tool generates proof obligations in an \emph{intermediate verification language}, from which prover-agnostic translators produce input for many provers at once.
Boogie~\cite{barnett2006boogie} pioneered this design and serves as the middle layer of, among others, the Dafny verifier~\cite{leino2010dafny}, which discharges its obligations with Z3.
Why3~\cite{bobot2011why3,filliatre2013why3} generalizes the idea: obligations are expressed in a polymorphic first-order logic and dispatched, through a system of \emph{drivers} that accommodate each prover's dialect and axiomatic strengths, to dozens of back-ends ranging from SMT solvers to saturation provers and proof assistants.
Industrial verifiers such as SPARK for Ada~\cite{hoang2015spark} and the WP plug-in of Frama-C for C~\cite{kirchner2015framac} rest on this infrastructure.
The B world joined this architecture when Mentr\'e et al.\ translated Atelier~B proof obligations into Why3's logic, modeling B's set theory as a polymorphic theory and measurably increasing the proportion of automatically discharged obligations on industrial case studies~\cite{mentre2012discharging}.

The pattern, however, concentrates trust in the middle: every driver, every encoding of the source theory, and every normalization performed by the platform belongs to the trusted base of the verification.
These translations are justified, when they are justified at all, by careful design and paper arguments; none of the platforms cited above verifies them mechanically.

\subsection{Encoding set theory for automated provers}
\label{subsec:rw-set-encodings}

The specification languages of the B family confront the translator with a distinctive gap.
Their expressions denote objects of a set-theoretic universe: membership is pervasive, functions are graphs, sets of sets are commonplace, and operators may be partial.
SMT solvers, by contrast, reason in many-sorted first-order logic over total functions, with no primitive notion of membership.
Every encoding must therefore decide how to represent sets---by an axiomatized membership predicate or by characteristic functions---and how to recover sorts from a language whose typing discipline, when it exists at all, is coarser than the solver's.
The systems reviewed below explore essentially every corner of this design space.

\paragraph*{Event-B and Rodin.}
D\'eharbe first isolated a class of Rodin proof obligations---over Booleans, integers, and basic sets---amenable to a direct SMT translation~\cite{deharbe2010automatic}.
The resulting SMT plug-in for Rodin~\cite{deharbe2012smt} industrialized two complementary encodings: a \emph{lambda-based} approach, which represents sets by their characteristic functions using the macro mechanism of SMT-LIB, and the \emph{ppTrans} approach~\cite{deharbe2014integrating}, which reuses the translator of Rodin's predicate prover to reduce set-theoretic constructs to first-order formulas over simple membership predicates.
On a large sample of industrial and academic projects, adding SMT solvers to the Atelier~B provers reduced the number of sequents requiring human interaction to roughly one fourth.

\paragraph*{Atelier~B and BWare.}
On the classical B side, the BWare project~\cite{delahaye2014bware} built a proof platform routing Atelier~B obligations, through Why3, to SMT solvers and first-order provers, with the Why3 encoding of Mentr\'e et al.~\cite{mentre2012discharging} as its front door.
A distinctive requirement of BWare was that provers deliver \emph{proof objects} that can be checked independently.
This requirement was met by Zenon Modulo~\cite{delahaye2013zenon}, an extension of the tableau prover Zenon~\cite{bonichon2007zenon} with typed proof search and deduction modulo, in which the axioms of B's set theory are turned into rewrite rules rather than clauses~\cite{bury2015automated}; the prover outputs certificates for the Dedukti universal proof checker, so that every proof found can be rechecked by a small trusted kernel~\cite{halmagrand2016automated}.
This line of work is the closest ancestor, in spirit, of the certification goal pursued in this thesis, with a crucial difference of scope discussed below: what is certified is each individual proof \emph{within} the axiomatized B theory, not the adequacy of that theory---or of the translation into it---with respect to B's semantics.

\paragraph*{ProB.}
The ProB toolset~\cite{leuschel2008prob} approaches B from the complementary direction of animation, model checking, and constraint solving, where the object is not to prove obligations but to find models and counterexamples.
The encoding questions are nevertheless the same.
ProB's kernel has been complemented by SAT-based model finding~\cite{plagge2012validating} and by two generations of Z3 integrations: first through a translation covering set comprehensions and other higher-order constructs, the two solvers exchanging intermediate results~\cite{krings2016smt}; then through a constructive lambda-based translation together with a CDCL(\(T\))-style solver implemented directly over ProB's own constraint kernel~\cite{schmidt2022smt}.
B2SAT~\cite{leuschel2024b2sat} completes the picture with a direct, low-level reduction of finite B constraint problems to propositional satisfiability.

\paragraph*{\tlaplus.}
The encoding problem has been studied with particular depth for \tlaplus~\cite{lamport2002specifying}, whose mathematical language combines \emph{untyped} Zermelo--Fraenkel set theory with choice, functions as primitive objects rather than sets of pairs, and arithmetic, underneath the Temporal Logic of Actions.
The \tlaplus Proof System (TLAPS)~\cite{cousineau2012tlaproofs} dispatches obligations to saturation provers, to an Isabelle encoding, and to SMT solvers.
Merz and Vanzetto developed the SMT backend through a series of encodings: a typed translation guided by type inference~\cite{merz2012smt}, then untyped and many-sorted encodings, injecting the set-theoretic universe into a single sort while recovering sorted reasoning where possible~\cite{merz2018encoding}, with type synthesis recast as constraint solving over refinement types in Vanzetto's thesis~\cite{vanzetto2014proof}.
The backend was recently redesigned by Defourn\'e with soundness as the explicit priority: the type-synthesis machinery, judged too complex to be trusted, was replaced by a direct first-order axiomatization of \tlaplus's set theory equipped with carefully chosen \emph{triggers} guiding quantifier instantiation, at essentially no cost in automation~\cite{defourne2023encoding}.
That redesign is a telling data point for this thesis: even within a single mature system, the translation itself---not the solver---came to be regarded as the weakest link.

\paragraph*{Kindred formalisms.}
Alloy~\cite{jackson2012software} restricts itself to a first-order relational logic precisely so that its analyzer can reduce verification questions to SAT within finite bounds~\cite{torlak2007kodkod}; El~Ghazi and Taghdiri later lifted Alloy assertions to unbounded proofs via SMT~\cite{elghazi2011relational}.
For Z, whose mathematical language closely prefigures B's, the Z/EVES system~\cite{saaltink1997zeves} combined rewriting and decision procedures long before the SMT era.

Across all these systems the pattern is uniform: the encodings embody substantial logical insight, their soundness is argued---sometimes meticulously---on paper, and the translators that implement them are complex, evolving, and unverified programs standing between the specification and the prover.

\subsection{Trusting the automation: certificates, reconstruction, and hammers}
\label{subsec:rw-certificates}

The question of trusting automated provers was posed early and sharply within interactive theorem proving, where a foreign proof must not compromise the assistant's kernel.
Barendregt and Barendsen distinguished the \emph{autarkic} answer---verify the external tool once and for all---from the \emph{skeptical} one---have it produce, on every run, a certificate that a small checker validates~\cite{barendregt2002autarkic}, the latter being an instance of the general paradigm of certifying algorithms~\cite{mcconnell2011certifying}.

The skeptical approach dominates practice.
Isabelle's Sledgehammer~\cite{paulson2010sledgehammer} translates goals to first-order provers and, since its extension to SMT solvers~\cite{blanchette2013extending}, reconstructs their answers inside the kernel~\cite{boehme2010z3,barbosa2020scalable}.
The proof format of veriT has since been standardized as Alethe~\cite{schurr2021alethe}, now emitted by cvc5 as well, checkable independently by the Carcara checker and elaborator~\cite{andreotti2023carcara}, and reconstructible in the Lambdapi proof assistant~\cite{coltellacci2026reconstruction}; cvc5 also exports proofs in other formats, including LFSC~\cite{stump2013lfsc} and Lean, at configurable granularity~\cite{barbosa2022flexible}.
In the Rocq ecosystem, SMTCoq validates SAT and SMT certificates through a checker programmed and \emph{proved correct} inside Rocq itself~\cite{armand2011modular,ekici2017smtcoq}; hammers in the same spirit exist for Rocq and HOL Light~\cite{czajka2018coqhammer,kaliszyk2014flyspeck}.
Dedukti~\cite{boespflug2012lambdapi,assaf2023dedukti} pushes the skeptical logic to its limit as a universal checker in which the proofs of many systems---including, as seen above, Zenon Modulo's proofs of B obligations---can be re-expressed and revalidated.
The autarkic route has been demonstrated as well, with a fully verified SAT solver with watched literals extracted from Isabelle~\cite{fleury2018sat}, though verifying a full SMT solver of competitive strength remains out of reach.

The Lean~4 ecosystem, in which this thesis is developed, has recently acquired the same infrastructure: Duper, a proof-producing superposition prover native to Lean~\cite{clune2024duper}; Lean-auto, a sound translation from Lean's dependent type theory to automated provers~\cite{qian2025leanauto}; lean-smt, which runs cvc5 and reconstructs its Alethe proofs as Lean terms~\cite{mohamed2025leansmt}; and bv\_decide, a bit-blasting decision procedure whose implementation is itself verified in Lean~\cite{boeving2025bvdecide}.

One structural remark matters for what follows.
In a hammer, the translation from the assistant's logic to the prover's need not be trusted: an unsound encoding can at worst produce an alien proof that fails to replay, never a false theorem.
The pipelines of \Cref{subsec:rw-set-encodings} enjoy no such safety net.
When Atelier~B or Rodin accepts a solver's \emph{unsat} verdict, the obligation is marked discharged; if the encoding that produced the SMT problem was unsound, a false obligation is silently certified, and no downstream certificate checking---which scrutinizes the solver's reasoning, not the translation---can detect it.
In these pipelines, the encoding is not merely part of the trusted base; it is its most exposed component.

\subsection{Certifying the translation itself}
\label{subsec:rw-certified-translation}

The remaining question---how to establish the correctness of the translation, rather than of the prover---has a classical analogue in compilation.
Translation validation checks each individual run of a translator a posteriori~\cite{pnueli1998translation}; verified translation proves the translator correct once and for all, as CompCert did for a realistic C compiler by programming and verifying it in Rocq~\cite{leroy2009compcert}.
Deductive verification has largely ignored both options for its own translators: the encodings of F\(^\star\) into Z3~\cite{fstar}, of Why3 to its many back-ends, and of the B-family tools surveyed above are all trusted.
Even the redesign of the TLAPS encoding undertaken with soundness as its explicit priority~\cite{defourne2023encoding} remains an informal argument about a nontrivial implementation, in which subtle defects have since been uncovered.

This manuscript takes the verified-translation route, applied for the first time, to our knowledge, to an SMT encoding of a set-theoretic specification language.
The encoder \beer{} is an executable Lean~4 program whose soundness is a theorem: source and target are interpreted in a single universe of ZFC sets, and the encoding provably preserves denotations~\cite{trelat2025beer}, as developed in \Cref{ch:zflean,ch:b,ch:smt,ch:loosening,ch:beer}.
For the translation, this is the autarkic route: the encoder is verified once and for all.
It combines naturally with the skeptical treatment of the solver, whose verdict can be validated by reconstructing its certificate in Lean in the style of lean-smt (\Cref{ch:barrel}), closing the chain from proof obligation to checked proof.
